\documentclass[10pt,a4paper]{amsart}
\usepackage[margin=19mm]{geometry}
\usepackage{amsmath,amssymb,mathtools}
\usepackage[scr=rsfs,cal=euler]{mathalfa}
\usepackage{xfrac}
\usepackage[unicode,hidelinks,bookmarks=false]{hyperref}
\newcommand{\ie}{i.e.\ }
\newcommand{\cf}{cf.\ }
\newcommand{\resp}{respectively\ }
\newcommand{\Subsection}{\subsection}
\providecommand{\nobreakdash}{\nobreak\hbox{-}}
\setlength{\emergencystretch}{2em}
\multlinegap0pt
\usepackage{bm}
\usepackage{bbm}
\usepackage{dsfont}
\usepackage{xspace}
\usepackage{cancel}
\usepackage{mathtools}
\usepackage{amsthm}
\makeatletter
\@ifundefined{theorem}{\newtheorem{theorem}{Theorem}}{}
\@ifundefined{lemma}{\newtheorem{lemma}[theorem]{Lemma}}{}
\makeatother
\def \loopmeasure{\mu^{\rm loop}}
\newcommand{\D}{\mathbb{D}}
\newcommand{\Lc}{\mathcal{L}}
\newcommand{\Prob}[1]{\mathbb{P} \left( #1 \right) }
\newcommand{\eps}{\varepsilon}
\renewcommand{\d}{\mathrm{d}}
\title[Crossing exponent in the Brownian loop soup]{Crossing exponent in the Brownian loop soup}
\author{Antoine Jego, Titus Lupu, Wei Qian}
\date{Source: arXiv:2303.03782v2 (CC BY 4.0)}
\begin{document}
\maketitle

\noindent\emph{Source and licence.} Crossing exponent in the Brownian loop soup. Version arXiv:2303.03782v2 (CC BY 4.0), distributed under the Creative Commons Attribution 4.0 International licence, as recorded in the arXiv licence metadata for this exact version and shown on the abstract page https://arxiv.org/abs/2303.03782v2. Full rights notice: licenses/arxiv-2303.03782-CC-BY-4.0.txt.\par\smallskip

\noindent\textit{Editorial scope.} Counted unit: the Lemma whose source label is \texttt{L:crossing\_equation\_sub} in the article (source0.tex lines 1318--1328, source label \texttt{L:crossing\_equation\_sub}), delivered together with the article's own complete original proof of it (source0.tex lines 1331--1441, one \texttt{proof} environment of 111 lines). No mathematics is written, repaired or re-derived: statement and proof are the article's own text line for line. Required context retained uncounted: L:loopmeasure\_annulus\_quantitative (lines 584--596); E:L\_loopmeasure\_annulus\_quantitative1 (lines 586--589); E:L\_loopmeasure\_annulus\_quantitative2 (lines 591--594); L:surround (lines 859--864). Every definition, display and label of those lines is retained unchanged. Omitted from this package and not redistributed: 1--583, 590--590, 595--858, 865--1317, 1329--1330, 1442--2657. Nothing inside a retained excerpt is omitted or altered: every retained body is delivered exactly as the article's lines are written, so no omission receipt of any kind accompanies this package. Citations: the retained excerpts carry no citation command at all, so no bibliography entry of the article is transcribed and no reference is redistributed. Reference adaptation: none was needed and none was made; every reference inside the delivered unit resolves inside it, so no unresolved reference or citation is produced. Printed slips kept literal: no printed word, symbol, hypothesis or number of the counted statement or of its proof was altered. Layout and class: the article's own class is \texttt{article} with packages including inputenc, fontenc, textcomp, amsmath, amssymb, amsthm, lmodern, geometry, xcolor, pict2e, microtype, listings, moreverb, hyperref; the wrapper is the lane's pinned \texttt{amsart} editorial wrapper at 10pt on A4, which restates the theorem-like environments the retained text needs and adds the article's own macro definitions that the retained text needs (\texttt{\textbackslash D}, \texttt{\textbackslash Lc}, \texttt{\textbackslash Prob}, \texttt{\textbackslash d}, \texttt{\textbackslash eps}, \texttt{\textbackslash loopmeasure}), with extra wrapper packages (bm, bbm, dsfont, xspace, cancel, mathtools). The delivered numbering is the unit's own. Assets: no retained span contains a figure environment, a graphics-inclusion command, a PDF-page inclusion, a table or a TikZ picture; no image file is copied into this package, the package declares no asset dependency and no image file is redistributed with it. Licence: version arXiv:2303.03782v2 of this article is distributed under the Creative Commons Attribution 4.0 International licence, as recorded in the arXiv licence metadata for this exact version and shown on the abstract page https://arxiv.org/abs/2303.03782v2. A full rights notice is delivered at licenses/arxiv-2303.03782-CC-BY-4.0.txt. Characters: every retained excerpt is pure ASCII, so the delivered engine, XeLaTeX with its default Latin Modern text font, reports no missing character.
\begin{lemma}\label{L:loopmeasure_annulus_quantitative}
Let $0<r_1<r_2<1$. We have
\begin{equation}
\label{E:L_loopmeasure_annulus_quantitative1}
\loopmeasure_\D ( \{ \wp \mathrm{~crossing~} r_2 \D \setminus r_1 \D \} ) = \left( 1 + O(1)\frac{r_1}{r_2} \right) \log \frac{\log (1/r_1)}{\log (r_2/r_1)}.
\end{equation}
Moreover, the left hand side of \eqref{E:L_loopmeasure_annulus_quantitative1} is differentiable with respect to $r_1$ and
\begin{equation}
\label{E:L_loopmeasure_annulus_quantitative2}
\frac{\d}{\d r_1} \loopmeasure_\D ( \{ \wp \mathrm{~crossing~} r_2 \D \setminus r_1 \D \} ) = \Big( 1 + \frac{O(1)}{\log(r_2/r_1)} \Big) \frac{1}{r_1} \Big( \frac{1}{\log(r_2/r_1)} - \frac{1}{\log(1/r_1)} \Big).
\end{equation}
In \eqref{E:L_loopmeasure_annulus_quantitative1} and \eqref{E:L_loopmeasure_annulus_quantitative2}, $O(1)$ refers to a quantity uniformly bounded with respect to $r_1$ and $r_2$.
\end{lemma}

\begin{lemma}\label{L:surround}
There exists $c = c(\theta)>0$ such that for all $r \in (0,1/10)$,
\begin{equation}
\Prob{ \exists \wp \in \Lc_\D^\theta \mathrm{~surrounding~} r \D } \geq 1 - r^c.
\end{equation}
\end{lemma}

\begin{lemma}\label{L:crossing_equation_sub}
Let $u \in (0,1)$. For any $\delta \in (0,1)$ and $s \geq 1$,
\begin{align*}
p_{\delta,\delta^s}
& = 1 - \left( \frac{s-1+u}{s} \right)^\theta + \frac{O(1)}{s \log \delta}
+ \left( u^\theta + \frac{O(1)}{\log \delta} \right) p_{\delta^u, \delta^{s-1+u}} + \left( 1 + O(1)\frac{(\log |\log \delta|)^C}{|\log \delta|^\theta} \right) \times \\
& \times \theta (1-u) (s-1+u)^\theta
\int_1^{1 + \frac{s-1}{u}} (s + (1-u)(t-1))^{-\theta-1} p_{(\delta^{s-1+u})^{1/t},\delta^{s-1+u}} \d t.
\end{align*}
The $O(1)$ appearing above stands for a quantity uniformly bounded with respect to $\delta$ and $s$ (it may depend on $u$).
\end{lemma}

\begin{proof}[Proof of Lemma \ref{L:crossing_equation_sub}]
Let $u \in (0,1)$, $s>1$ and $\delta >0$. By scaling,
\[
p_{\delta,\delta^s} = \Prob{ \delta^u \overset{\Lc_{\delta^{-1+u} \D}}{\longleftrightarrow} \delta^{s-1+u} }.
\]
Let
\begin{equation}
t_* = \min \{ t \geq 1: \exists \wp \in \Lc_{ \delta^{-1+u} \D}^\theta \text{~crossing~} \D \setminus ( \delta^{s-1+u})^{1/t} \D \}.
\end{equation}
If $t_* =1$, then it means that the crossing that we are interested in is entirely made by a single loop which visits $\D^c$. If $t_* >1$, then we need loops included in $\D$ to finish the crossing. Notice further that if $t \geq 1 + \frac{s-1}{u}$, then $(\delta^{s-1+u})^{1/t}$ is larger than $\delta^u$ and the crossing of the annulus $\delta^u \D \setminus \delta^{s-1+u} \D$ is entirely made by loops staying in $\D$. Putting these remarks together, we obtain that
\begin{align*}
p_{\delta,\delta^s}
& = \Prob{ t_* = 1} + \Prob{t_* \geq 1 + \frac{s-1}{u} } p_{\delta^u, \delta^{s-1+u}} \\
& + \int_1^{1 + \frac{s-1}{u}} \Prob{ t_* \in \d t}
\Prob{ \delta^{u} \overset{\Lc_{\delta^{-1+u} \D}^\theta}{\longleftrightarrow} \delta^{s-1+u} \Big\vert t_* = t}.
\end{align*}
We are now going to estimate precisely each term appearing in the above display. These terms are divided into two different types of probabilities: probabilities that either a \emph{loop} or a \emph{cluster} crosses some annulus. Terms of the first type can be explicitly computed thanks to Lemma \ref{L:loopmeasure_annulus_quantitative}, whereas terms of the second type will expressed in terms of $F_\delta$.
In the sequel, the $O(1)$ terms might depend on $u$. In particular, the fact that $u$ is positive will be important for the estimates we are going to write.

By scaling and by Lemma \ref{L:loopmeasure_annulus_quantitative},
\begin{align*}
& \Prob{t_* = 1} =  \Prob{ \exists \wp \in \Lc_{\delta^{-1+u} \D}^\theta \text{~crossing~} \D \setminus \delta^{s-1+u} \D}
= \Prob{ \exists \wp \in \Lc_{\D}^\theta \text{~crossing~} \delta^{1-u} \D \setminus \delta^s \D} \\
& = 1 - \exp \left( - \theta \loopmeasure_\D (\{\wp \in \Lc_{\D}^\theta \text{~crossing~} \delta^{1-u} \D \setminus \delta^s \D\}) \right)
= 1 - \left( \frac{s-1+u}{s} \right)^\theta + \frac{O(1)}{s \log \delta}.
\end{align*}
Similarly,
\begin{align*}
\Prob{t_* \geq 1 + \frac{s-1}{u} } = \Prob{ \nexists \wp \in \Lc_{\delta^{-1+u}}^\theta \text{~crossing~} \D \setminus \delta^u \D }
= u^\theta + \frac{O(1)}{\log \delta}.
\end{align*}
Finally, writing $r_t = \delta^{1-u + \frac{s-1+u}{t}}$,
\begin{align*}
& \Prob{t_* \in \d t} = \frac{\d}{\d t} \Prob{ \exists \wp \in \Lc_{\delta^{-1+u} \D}^\theta \text{~crossing~} \D \setminus (\delta^{s-1+u})^{1/t} \D} \d t \\
& = \theta \left( \frac{\d}{\d t} \loopmeasure_\D ( \{ \wp \in \Lc_{\D}^\theta \text{~crossing~} \delta^{1-u} \D \setminus r_t \D \} ) \right) e^{ - \theta \loopmeasure_\D ( \{ \wp \in \Lc_{\D}^\theta \text{~crossing~} \delta^{1-u} \D \setminus r_t \D \} ) } \d t.
\end{align*}
By Lemma \ref{L:loopmeasure_annulus_quantitative},
\begin{align*}
& \frac{\d}{\d t} \loopmeasure_\D ( \{ \wp \in \Lc_{\D}^\theta \text{~crossing~} \delta^{1-u} \D \setminus r_t \D \} ) = \left( \frac{\d r_t}{\d t} \right) \frac{\d}{\d r} \loopmeasure_\D ( \{ \wp \in \Lc_{\D}^\theta \text{~crossing~} \delta^{1-u} \D \setminus r \D \} ) \Big\vert_{r = r_t} \\
& = \left( 1 + \frac{O(1)}{\frac{s-1+u}{t} \log \delta} \right)  \left( \frac{s-1+u}{t^2} \log (\delta) r_t \right) \frac{1}{r_t} \frac{1}{\log \delta} \left( \frac{1}{\frac{s-1+u}{t}} - \frac{1}{1-u + \frac{s-1+u}{t}} \right) \\
& = \left( 1 + \frac{O(1)}{\log \delta} \right) \frac{1-u}{t(1-u)+s-1+u}.
\end{align*}
Together with \eqref{E:L_loopmeasure_annulus_quantitative1}, it gives that
\begin{align*}
& \Prob{t_* \in \d t}
= \left( 1 + \frac{O(1)}{\log \delta} \right) \theta \frac{1-u}{t(1-u)+s-1+u}
\left( \frac{(1-u)t + s-1+u}{s-1+u} \right)^{-\theta} \d t \\
& = \left( 1 + \frac{O(1)}{\log \delta} \right) \theta (1-u) (s-1+u)^\theta (s + (1-u)(t-1))^{-\theta-1} \d t.
\end{align*}
Putting things together, we have obtained that
\begin{align*}
p_{\delta,\delta^s}
& = 1 - \left( \frac{s-1+u}{s} \right)^\theta + \frac{O(1)}{s \log \delta}
+ \left( u^\theta + \frac{O(1)}{\log \delta} \right) p_{\delta^u, \delta^{s-1+u}} + \left( 1 + \frac{O(1)}{\log \delta} \right) \times \\
& \times \theta (1-u) (s-1+u)^\theta
\int_1^{1 + \frac{s-1}{u}} (s + (1-u)(t-1))^{-\theta-1} \Prob{ \delta^u \overset{\Lc_{\delta^{-1+u} \D}^\theta}{\longleftrightarrow} \delta^{s-1+u} \Big\vert t_* = t} \d t.
\end{align*}

It remains to estimate the conditional probability appearing in the integral above.
Let $t \in [1, 1 + \frac{s-1}{u}]$.
We are going to conclude the proof by showing that
\begin{equation}
\label{E:proof_1234}
\Prob{ \delta^u \overset{\Lc_{\delta^{-1+u} \D}^\theta}{\longleftrightarrow} \delta^{s-1+u} \Big\vert t_* = t}
= \left( 1 - O(1) \frac{\log |\log \delta|^C}{|\log \delta|^\theta} \right) p_{(\delta^{s-1+u})^{1/t},\delta^{s-1+u}}.
\end{equation}
The upper bound is clear. Indeed,
conditionally on $t_* = t$, to make the crossing from $\delta^{s-1+u}$ to $\delta^u$, the loops which stay in $\D$ have to cross at least the annulus $(\delta^{s-1+u})^{1/t} \D \setminus \delta^{s-1+u} \D$. This is saying that
\[
\Prob{ \delta^u \overset{\Lc_{\delta^{-1+u} \D}^\theta}{\longleftrightarrow} \delta^{s-1+u} \Big\vert t_* = t}
\leq p_{(\delta^{s-1+u})^{1/t}, \delta^{s-1+u}}.
\]
We will now provide a lower bound. The idea is to ``weld'' the crossing from $\delta^u \D$ to $(\delta^{s-1+u})^{1/t} \D$, and the crossing from $(\delta^{s-1+u})^{1/t} \D$ to $(\delta^{s-1+u}) \D$. Let $t_+$ and $t_-$ be close to $t$ and be such that $t_- < t < t_+$.
Conditionally on $t_* = t$, to create a crossing of the annulus $\delta^u \D \setminus \delta^{s-1+u} \D$, it is enough for the loops staying in $\D$ to do the following four events:
\begin{itemize}
\item
Event $E_1$: There is a cluster crossing the annulus $(\delta^{s-1+u})^{1/t} \D \setminus \delta^{s-1+u} \D$;
\item
Event $E_2$: There is a loop staying in $(\delta^{s-1+u})^{1/t} \D$ which surrounds $(\delta^{s-1+u})^{1/t_-} \D$;
\item
Event $E_3$: There is a loop staying in $(\delta^{s-1+u})^{1/t_+} \D$ which surrounds $(\delta^{s-1+u})^{1/t} \D$;
\item
Event $E_4$: There is a loop crossing the annulus $(\delta^{s-1+u})^{1/t_+} \D \setminus (\delta^{s-1+u})^{1/t_-} \D$
\end{itemize}
By FKG inequality, the probability of the intersection of these four events is at least the product of each probability. This gives
\begin{align*}
& \Prob{ \delta^{s-1+u} \overset{\Lc_{\delta^{-1+u} \D}^\theta}{\longleftrightarrow} \delta^u \Big\vert t_* = t}
\geq p_{(\delta^{s-1+u})^{1/t},\delta^{s-1+u}} \Prob{E_2} \Prob{E_3} \Prob{E_4}.
\end{align*}
Let us denote $\eps = \delta^{s-1+u}$.
By Lemma \ref{L:surround}, there exists $c>0$ such that
$
\Prob{ E_2 } \geq 1 - \eps^{c (\frac{1}{t_-} - \frac1t)}.
$
We choose $t_+$ and $t_-$ precisely so that
\[
\frac{1}{t_-} - \frac{1}{t} = \frac1t - \frac{1}{t_+} = \frac{\theta}{c} \frac{\log \log \frac1\eps}{\log \frac1\eps}.
\]
This choice leads to:
$\Prob{ E_2 } = \Prob{E_3} \geq 1 - |\log \eps|^{-\theta}.$
On the other hand, by Lemma \ref{L:loopmeasure_annulus_quantitative},
\[
\Prob{ E_4 } = 1 - \left( \frac{\frac{1}{t_-} - \frac{1}{t_+}}{ \frac{1}{t_-} } \right)^{\theta + o(1/\log \delta)} = 1 - O(1) t^{\theta} |\log \eps|^{-\theta} \log |\log \delta|^C = 1 - O(1) \frac{\log |\log \delta|^C}{|\log \delta|^\theta}.
\]
Overall, we have obtained 
\[
\Prob{ \delta^u \overset{\Lc_{\delta^{-1+u} \D}^\theta}{\longleftrightarrow} \delta^{s-1+u} \Big\vert t_* = t}
\geq \left( 1 - O(1) \frac{\log |\log \delta|^C}{|\log \delta|^\theta} \right) p_{(\delta^{s-1+u})^{1/t},\delta^{s-1+u}}.
\]
This concludes the proof of \eqref{E:proof_1234} and the proof of Lemma \ref{L:crossing_equation_sub}.
\end{proof}

\end{document}
